\title{DeepSeekMath: Pushing the Limits of Mathematical Reasoning in Open Language Models}

\begin{abstract}
The task of mathematical reasoning remains particularly challenging for language models, given the intricate and systematic features inherent to mathematics. In this study, we present DeepSeekMath~7B, an extension of DeepSeek-Coder-Base-v1.5 7B through continued pre-training using 120B math-specific tokens extracted from Common Crawl, alongside supplemental natural language and code corpora. DeepSeekMath~7B demonstrates a remarkable performance, securing a 51.7\% score on the high-level MATH benchmark, all without the need for supplementary toolkits or ensemble strategies, thus nearing the results achieved by Gemini-Ultra and GPT-4. When evaluating self-consistency over 64 outputs, DeepSeekMath~7B attains 60.9\% on MATH. The superior mathematical reasoning exhibited by DeepSeekMath~can be attributed to two principal elements: firstly, our utilization of a carefully constructed data selection pipeline to exploit the full capacity of publicly accessible web resources; secondly, our development of Group Relative Policy Optimization (GRPO), an adaptation of Proximal Policy Optimization (PPO), which both augments mathematical problem-solving proficiency and streamlines the memory efficiency of PPO.
\end{abstract}

\section{Introduction}

Large language models (LLMs) have transformed methodologies for addressing mathematical reasoning within artificial intelligence, catalyzing progress on quantitative benchmarks~\citep{MATH} and on geometric reasoning evaluations~\citep{trinh2024solving}. Beyond automated tasks, these models increasingly support humans in tackling advanced mathematical challenges~\citep{tao}. Nonetheless, state-of-the-art systems such as GPT-4~\citep{gpt4} and Gemini-Ultra~\citep{gemini} remain closed-source, and existing open-access models lag significantly in their mathematical capabilities.

This work presents DeepSeekMath, a specialized language model focused on mathematical tasks that surpasses existing open-source alternatives and closely rivals the performance of GPT-4 across a range of academic benchmarks. Central to our approach is the construction of the DeepSeekMath Corpus, a massive pre-training dataset containing 120B high-quality math tokens. To compile this dataset, we leverage a fastText-based classifier \citep{joulin2016fasttext} to filter content from the Common Crawl (CC). Initially, this classifier is trained using examples from OpenWebMath \citep{openwebmath} as positive samples, with a varied collection of unrelated web pages used as negative samples. The trained classifier is then applied to the CC to identify further positive samples, which are subsequently validated and refined through human annotation. These newly labeled data augment the training set, leading to iterative improvements of the classifier. Empirical results demonstrate that the corpus maintains a high standard, as evidenced by our DeepSeekMath-Base 7B model achieving a 64.2\% accuracy on GSM8K \citep{gsm8k} and 36.2\% on the advanced MATH benchmark \citep{MATH}, both figures exceeding the results of Minerva 540B \citep{minerva}. Furthermore, since the DeepSeekMath Corpus includes multilingual content, our model also shows enhanced performance on Chinese math benchmarks \citep{wei2023cmath,agieval}. We contend that our methodology for mathematical data curation offers a valuable reference for ongoing research, and we anticipate considerable advancements in the near future.

DeepSeekMath-Base is initialized using DeepSeek-Coder-Base-v1.5 7B \citep{deepseek-coder}, as our findings indicate that leveraging a model pre-trained on code yields better outcomes than starting with a generic language model. Moreover, we find that the process of training on mathematical data not only elevates mathematical competence but also enhances performance on general reasoning benchmarks, such as MMLU \citep{mmlu} and BBH \citep{bbh}.

Following pre-training, we further refine DeepSeekMath-Base through mathematical instruction tuning, incorporating data that supports chain-of-thought reasoning \citep{cot}, program-of-thought execution \citep{pot,pal}, and tool-augmented reasoning \citep{tora}. The resulting model, DeepSeekMath-Instruct 7B, surpasses all other 7B-sized models and reaches a performance level on par with 70B parameter open-source instruction-tuned models.

Additionally, we propose Group Relative Policy Optimization (GRPO), a reinforcement learning algorithm derived from Proximal Policy Optimization (PPO) \citep{schulman2017proximal}. Unlike traditional PPO, GRPO dispenses with the critic model, instead deriving the baseline directly from group-based reward aggregation, leading to considerable savings in training computation. Employing only a selected portion of English instruction tuning data, GRPO drives marked improvements over the already strong DeepSeekMath-Instruct model, across both within-domain tasks (GSM8K: 82.9\% $\rightarrow$ 88.2\%, MATH: 46.8\% $\rightarrow$ 51.7\%) and out-of-domain mathematical benchmarks (e.g., CMATH: 84.6\% $\rightarrow$ 88.8\%) in the reinforcement learning phase. 

We also introduce a unified conceptual framework to interpret a range of methodologies—such as Rejection Sampling Fine-Tuning (RFT) \citep{yuan2023scaling}, Direct Preference Optimization (DPO) \citep{dpo}, PPO, and GRPO—demonstrating that all of them are fundamentally rooted in either direct or streamlined forms of reinforcement learning. Through comprehensive experiments—including analyses of online versus offline training, outcome-based versus process-based supervision, and single-turn versus iterative RL—we thoroughly investigate the core components underlying this paradigm. Finally, we discuss the mechanisms behind the reinforcement learning enhancements seen in instruction-tuned models and outline future research avenues for advancing RL effectiveness within this unified framework.

\subsection{Contributions}
We present advancements in large-scale mathematical pre-training and a thorough investigation of reinforcement learning techniques.

\textbf{Math Pre-Training at Scale}

\begin{itemize}[topsep=0pt]
    \item Our study demonstrates that Common Crawl, a publicly available resource, harbors substantial mathematically-relevant data. Through the deployment of a rigorously curated data selection mechanism, we assembled the DeepSeekMath Corpus—a high-quality dataset comprising 120B tokens sourced from mathematics-oriented web content. This corpus surpasses the scale of the datasets employed by Minerva \citep{minerva} by nearly sevenfold, and exceeds that of OpenWebMath \citep{openwebmath} by approximately ninefold.

    \item The DeepSeekMath-Base 7B pre-trained model we developed exhibits performance that rivals Minerva 540B \citep{minerva}, suggesting that model size is not the sole determinant of proficiency in mathematical reasoning. Our results show that smaller-scale models, when trained on curated and high-caliber data, can yield strong performance.

    \item We document insights from our experiments on math-focused model training. Pre-training on code, before engaging in mathematical instruction, enhances the model’s capability to solve math problems regardless of tool utilization. This sheds light on a longstanding topic: \textit{does code training enhance reasoning skills?} Our findings support this for the mathematical domain.

    \item Contrary to widespread practice, augmenting training sets with arXiv papers—an approach prevalent in math-oriented works—fails to provide appreciable improvements across the mathematical benchmarks considered in this study.
\end{itemize}

\textbf{Exploration and Analysis of Reinforcement Learning}

\begin{itemize}[topsep=0pt]
    \item This work presents Group Relative Policy Optimization (GRPO), a reinforcement learning technique that is both efficient and effective. Unlike traditional approaches that employ a critic network, GRPO computes the baseline using group-based scoring, thus considerably lowering computational demands relative to Proximal Policy Optimization (PPO).
    \item We show that employing GRPO yields marked improvements in the capabilities of our instruction-tuned model, DeepSeekMath-Instruct, when utilizing only instruction-tuning datasets. Additionally, reinforcement learning with GRPO leads to observed gains in the model’s ability to generalize to domains outside the training distribution.
    \item A unified conceptual framework is introduced to compare and contextualize different methods such as RFT, DPO, PPO, and GRPO. We further perform thorough experimentation—including comparisons between online and offline learning, outcome-focused versus process-oriented supervision, as well as single-step versus iterative reinforcement learning—to systematically analyze the core factors underlying this framework.
    \item Building on this integrative framework, we examine the drivers behind reinforcement learning's empirical benefits and outline several prospective research avenues to foster more efficient reinforcement learning for large language models (LLMs).
\end{itemize}

\subsection{Summary of Evaluations and Metrics}

\begin{itemize}[topsep=0pt]
    \item \textbf{English and Chinese Mathematical Reasoning}: We carry out extensive model evaluations using both English and Chinese datasets, addressing mathematical questions ranging from elementary to collegiate complexity. English benchmarks encompass GSM8K \citep{gsm8k}, MATH \citep{MATH}, SAT \citep{llemma}, OCW Courses \citep{minerva}, and MMLU-STEM \citep{mmlu}. For Chinese tasks, MGSM-zh \citep{mgsm}, CMATH \citep{wei2023cmath}, Gaokao-MathCloze \citep{agieval}, and Gaokao-MathQA \citep{agieval} are used. We assess not only the models’ competence in producing comprehensive textual solutions unaided by external tools, but also their proficiency in solving problems via Python programming.

    On the English evaluation sets, DeepSeekMath-Base demonstrates parity with the proprietary Minerva 540B \citep{minerva} and consistently outperforms all open-source models (e.g., Mistral 7B \citep{mistral} and Llemma-34B \citep{llemma}), independent of whether they have benefited from math-specific pre-training, frequently by a notable gap. Remarkably, DeepSeekMath-Base also establishes a lead on Chinese evaluation sets, attributable in part to our inclusion of high-caliber non-English training data, contrasting with prior studies \citep{minerva,llemma} which predominantly utilized English-centric datasets. Through the combination of mathematical instruction fine-tuning and reinforcement learning, both DeepSeekMath-Instruct and DeepSeekMath-RL exhibit exceptional outcomes, surpassing 50\% accuracy on the highly challenging MATH benchmark—setting a new milestone for open-source models on this task.

\item \textbf{Formal Mathematics}: We assess DeepSeekMath-Base's capabilities using the informal-to-formal theorem proving benchmark introduced in \citep{dsp_proof}, applying it to the miniF2F dataset \citep{minif2f} and leveraging Isabelle \citep{isabelle} as the proof assistant. DeepSeekMath-Base achieves robust few-shot performance in the task of autoformalization.

\item \textbf{Natural Language Understanding, Reasoning, and Code}: To comprehensively evaluate the model’s aptitude in language understanding, logical reasoning, and code generation, DeepSeekMath-Base is tested on the Massive Multitask Language Understanding (MMLU) benchmark \citep{mmlu}, which spans 57 diverse multiple-choice topics. Additionally, it is evaluated on BIG-Bench Hard (BBH) \citep{bbh}, consisting of 23 particularly difficult multi-step reasoning tasks, as well as on HumanEval \citep{codex} and MBPP \citep{mbpp}, both standard datasets for appraising code language models. The results show that math pre-training significantly enhances both language comprehension and reasoning accuracy.

\section{Math Pre-Training}

\subsection{Data Collection and Decontamination}

This section details the method used to assemble the DeepSeekMath Corpus from the Common Crawl dataset. As illustrated in Figure \ref{fig:data_collect}, we introduce an iterative pipeline designed for large-scale extraction of mathematical content, initiating from a seed dataset—typically a smaller, high-quality selection of mathematics-related texts. This methodology is flexible and can be adapted to other fields, such as programming.

Initially, OpenWebMath \citep{openwebmath}, which comprises rigorously vetted mathematical web content, is selected as the seed dataset. A fastText model \citep{joulin2016fasttext} is then trained to help identify and retrieve additional web pages similar to those in OpenWebMath. Specifically, we sample 500,000 entries from the seed set as positive samples and select another 500,000 Common Crawl pages as negative samples. Model training is performed using an open-source library\footnote{\url{https://fasttext.cc}}, with hyperparameters set to a vector size of 256, a learning rate of 0.1, a word n-gram window up to 3, a minimum word occurrence threshold of 3, and 3 training epochs. 

To manage Common Crawl's scale, we apply deduplication techniques based on URL matching and content similarity, reducing the dataset to 40 billion HTML pages. The trained fastText model is then employed to retrieve candidate mathematical pages from this deduplicated set. Collected pages are subsequently ranked by their fastText model score, with only the highest-scoring pages being retained. We empirically evaluate data selection by conducting pre-training runs with datasets containing the top 40B, 80B, 120B, and 160B tokens, ultimately selecting the top 40B tokens during the first round.

Following the initial round of data collection, a significant portion of mathematical web pages are not retrieved. This shortfall primarily results from the limited diversity in the positive examples used to train the fastText model. To address this, we seek out further mathematical web sources to expand the seed dataset and enhance the fastText model's performance. The methodology involves first partitioning the full Common Crawl dataset into non-overlapping domains, where a domain comprises web pages sharing a base URL. For each domain, we determine the proportion of web pages that were acquired in the first iteration. Any domain with more than 10\% of its pages collected is labeled as math-related (e.g., \url{mathoverflow.net}).

We then proceed to manually annotate the subset of URLs within these math-related domains that correspond to mathematical content (for instance, \url{mathoverflow.net/questions}). Web pages associated with these annotated URLs, if not already collected, are incorporated into the seed dataset. This strategy ensures an increased supply of positive examples, leading to the training of a more effective fastText model that can retrieve a greater quantity of mathematical data in later rounds. After four iterations, the resulting collection contains 35.5M mathematical web pages, amassing 120B tokens. During the fourth iteration, it is observed that approximately 98\% of the data had already been captured in the third iteration, prompting a decision to stop further collection.

To mitigate the risk of contaminating evaluation benchmarks, we adhere to the filtering protocol of \cite{deepseek-coder}, excluding web pages that feature questions or answers from established English mathematical benchmarks like GSM8K \citep{gsm8k} and MATH \citep{MATH}, as well as Chinese datasets such as CMATH \citep{wei2023cmath} and AGIEval \citep{agieval}. The procedure specifies that any text segment containing an exact 10-gram string match with any portion of a benchmark dataset is removed from the math training corpus. For benchmark passages comprising fewer than 10 grams but at least 3 grams, we employ strict matching to exclude contaminated web pages.

\subsection{Validating the Quality of the DeepSeekMath~Corpus}

To assess the quality of the DeepSeekMath~Corpus, we conduct pre-training experiments in comparison to the latest mathematics-oriented corpora:

\begin{itemize}[topsep=0pt]
    \item \textbf{MathPile} \citep{mathpile}: a collection (8.9B tokens) assembled from multiple sources, including textbooks, Wikipedia, ProofWiki, CommonCrawl, StackExchange, and arXiv, with the latter constituting over 85\% of the data;
    \item \textbf{OpenWebMath} \citep{openwebmath}: filtered CommonCrawl content focused on mathematical material, containing a total of 13.6B tokens;
    \item \textbf{Proof-Pile-2} \citep{llemma}: a composite mathematics corpus combining OpenWebMath, AlgebraicStack (10.3B tokens of mathematical code), and arXiv papers (28.0B tokens). In our experiments on Proof-Pile-2, we adopt the arXiv:Web:Code token ratio of 2:4:1 as described in \cite{llemma}.
\end{itemize}

\subsubsection{Training Setting}
\label{sec:quality-policy}
We fine-tune a general-purpose language model with 1.3B parameters—architecturally consistent with DeepSeek LLMs \citep{deepseek-llm} and referenced here as DeepSeek-LLM 1.3B—on individual mathematical datasets, training each for 150B tokens. All model training utilizes the resource-efficient HAI-LLM framework \citep{haillm}. Following DeepSeek’s established methodology, we adopt the AdamW optimizer \citep{adamW} with parameters $\beta_1=0.9$, $\beta_2=0.95$, and $\mathrm{weight\_decay}=0.1$. The learning rate schedule consists of an initial 2,000-step linear warmup to peak value, followed by a decrease to 31.6\% at 80\% of training completion, and a further reduction to 10.0\% of peak at 90\% completion. The highest learning rate is capped at 5.3e-4. We use a batch size equivalent to 4M tokens, and each training instance processes up to 4K tokens in context.

\subsubsection{Evaluation Results}

\textbf{The DeepSeekMath~Corpus distinguishes itself by its high quality, comprehensive multilingual scope, and unmatched scale.}

\begin{itemize}[topsep=0pt]
    \item \textbf{High-quality}:
    To assess model capabilities, we conduct few-shot chain-of-thought prompting evaluations \cite{cot} across eight mathematical benchmarks. As detailed in Table \ref{tab:corpora_comparison}, DeepSeek-LLM 1.3B trained on DeepSeekMath~Corpus consistently achieves superior results. Figure \ref{fig:corpora_comparisons} further demonstrates its performance surpasses that of the Proof-Pile-2-trained model at 50B tokens (covering one epoch of Proof-Pile-2), signifying the overall higher data quality of DeepSeekMath~Corpus.
    \item \textbf{Multilingual}:
    The DeepSeekMath~Corpus contains mathematical texts in a variety of languages, with English and Chinese constituting the primary linguistic segments. Results in Table \ref{tab:corpora_comparison} indicate that pretraining on DeepSeekMath~Corpus enhances the model’s mathematical reasoning in both English and Chinese. Conversely, prevailing corpora with their English-centric compositions exhibit modest gains in English and negligible or even negative effects on Chinese-language mathematical reasoning.
    \item \textbf{Large-scale}:
    In terms of volume, the DeepSeekMath~Corpus dwarfs other existing mathematical datasets. Figure \ref{fig:corpora_comparisons} illustrates that DeepSeek-LLM 1.3B, when exposed to the DeepSeekMath~Corpus, not only achieves a steeper learning trajectory but also maintains gains over longer training spans. The much smaller baseline corpora are rapidly exhausted, leading to repeated dataset passes and early stagnation of model improvements.
\end{itemize}

\subsection{Training and Evaluating DeepSeekMath-Base 7B}

This section presents DeepSeekMath-Base 7B, a foundational model engineered for advanced mathematical reasoning. The model's parameters are initialized from DeepSeek-Coder-Base-v1.5 7B \citep{deepseek-coder}, and it undergoes training on a corpus totaling 500B tokens. The data composition consists of 56\% DeepSeekMath Corpus, 4\% AlgebraicStack, 10\% arXiv materials, 20\% code sourced from Github, and the final 10\% comprising natural language texts from Common Crawl in both English and Chinese. The training protocol largely aligns with the configuration detailed in Section \ref{sec:quality-policy}, with the exceptions that the learning rate peaks at 4.2e-4 and training proceeds with a batch size of 10M tokens.

We systematically evaluate DeepSeekMath-Base 7B’s proficiency in mathematical reasoning, examining its capability to independently generate complete mathematical solutions, leverage external tools for problem-solving, and engage in formal theorem verification. Furthermore, we report the model's broader competencies, detailing results in natural language understanding, reasoning, and programming tasks.

\paragraph{Mathematical Problem Solving with Step-by-Step Reasoning}

To assess DeepSeekMath-Base’s problem-solving abilities, we employ few-shot chain-of-thought prompting \citep{cot} across eight established benchmarks in both English and Chinese. The benchmarks address various aspects of quantitative reasoning, such as GSM8K \citep{gsm8k}, MATH \citep{MATH}, and CMATH \citep{wei2023cmath}, as well as multiple-choice evaluations like MMLU-STEM \citep{mmlu} and Gaokao-MathQA \citep{agieval}. This range captures mathematical problems spanning primary through university-level subjects, offering a comprehensive gauge of the model's capabilities.

As indicated in Table \ref{tab:base_math_cot}, among publicly available base models, DeepSeekMath-Base 7B achieves the highest scores on all eight evaluation benchmarks, surpassing popular models such as Mistral 7B \citep{mistral} and the newly introduced Llemma 34B \citep{llemma}, which was trained on the Proof-Pile-2 dataset \citep{llemma}. Particularly on the challenging MATH dataset, DeepSeekMath-Base 7B exceeds the top open-source competitors by more than 10 percentage points in absolute terms, and even surpasses Minerva 540B \citep{minerva}—a much larger, proprietary base model derived from PaLM \citep{palm} and specifically fine-tuned on mathematical corpora.

\paragraph{Mathematical Problem Solving with Tool Use} To assess program-assisted mathematical reasoning, we employ few-shot program-of-thought prompting \citep{pot,pal} on both GSM8K and MATH datasets. The model is tasked with constructing a Python script to address each problem, making use of computational packages like \textit{math} and \textit{sympy} for sophisticated calculations, with the answer determined by the program’s output. Table \ref{tab:base_math_pot_proof} demonstrates that DeepSeekMath-Base 7B sets a new standard, outperforming the previous best open-source model, Llemma 34B.

\paragraph{Formal Mathematics}

The automation of formal proofs is essential for ensuring correctness, reproducibility, and productivity in mathematical reasoning and has attracted growing research interest. We examine the performance of DeepSeekMath-Base 7B in informal-to-formal proof generation, as formulated by \citep{dsp_proof}: transforming an informal problem statement, its corresponding formal version, and a non-rigorous proof into a formal proof. Experiments are performed on the miniF2F benchmark \citep{minif2f}, which targets advanced formal mathematics typical of mathematical Olympiads, using few-shot prompting to create proofs in Isabelle for each challenge. In alignment with \cite{dsp_proof}, the model produces a structured proof outline that is then completed by invoking Sledgehammer \citep{sledgehammer}, an automated theorem prover. According to the results in Table \ref{tab:base_math_pot_proof}, DeepSeekMath-Base 7B displays robust capability in generating formalized proofs automatically.

\paragraph{Natural Language Understanding, Reasoning, and Code}

We further investigate model abilities on natural language comprehension via MMLU \citep{mmlu}, complex reasoning tasks via BBH \citep{bbh}, and programming skills with HumanEval \citep{codex} and MBPP \citep{mbpp}. Results presented in Table \ref{tab:understanding_reasoning_code} reveal that DeepSeekMath-Base 7B notably improves over its predecessor, DeepSeek-Coder-Base-v1.5 \citep{deepseek-coder}, especially on MMLU and BBH, highlighting the beneficial effects of mathematical training for both language understanding and reasoning. The inclusion of code-related tokens throughout continual training enables DeepSeekMath-Base 7B to match the coding benchmark scores previously set by DeepSeek-Coder-Base-v1.5. In summary, DeepSeekMath-Base 7B establishes a considerable lead over the general-purpose Mistral 7B \citep{mistral} across the assessed reasoning and code-generation tasks.

\section{Supervised Fine-Tuning}

\subsection{SFT Data Curation}

We assemble a dataset for mathematical instruction-tuning that incorporates both English and Chinese tasks drawn from a spectrum of mathematical disciplines and complexity levels. Each problem is aligned with a solution presented in one of several formats: chain-of-thought (CoT) \citep{cot}, program-of-thought (PoT) \citep{pot,pal}, or tool-integrated reasoning \citep{tora}. The resulting training corpus contains 776K examples.

\begin{itemize}[topsep=0pt]
    \item \textbf{English mathematical datasets}: Our annotation process applies tool-integrated reasoning to GSM8K and MATH problems, and additionally utilizes a selected subset of MathInstruct \citep{MathInstruct} together with the Lila-OOD \citep{lila} training set, both featuring solutions formulated as CoT or PoT. This English dataset spans a broad array of mathematical domains such as algebra, probability, number theory, calculus, and geometry.
    \item \textbf{Chinese mathematical datasets}: We gather K-12 mathematics questions in Chinese, covering 76 distinct sub-topics (e.g., linear equations), with accompanying solutions written in both CoT and tool-integrated reasoning styles.
\end{itemize}

\subsection{Training and Evaluating DeepSeekMath-Instruct 7B}

This section details the development of DeepSeekMath-Instruct 7B, which is refined through instruction tuning built upon DeepSeekMath-Base. To construct the input for each training step, problems are concatenated in a random order until the length approaches the 4K token context window. The model is optimized over 500 steps using a batch size of 256 and a fixed learning rate of 5e-5.

For evaluation, we assess mathematical reasoning proficiency of the models—with and without the use of external tools—across four quantitative reasoning benchmarks in both English and Chinese. The comparison includes contemporaneous state-of-the-art models:

\begin{itemize}[topsep=0pt]
    \item \textbf{Closed-source models}: Our baseline models comprise (1) the GPT series, notably GPT-4 \citep{gpt4} and GPT-4 Code Interpreter~\footnote{\url{https://openai.com/blog/chatgpt-plugins\#\#code-interpreter}}, (2) Gemini Ultra and Pro \citep{gemini}, (3) Inflection-2 \citep{inflection-2}, (4) Grok-1~\footnote{\url{https://x.ai/model-card}}, and models released by Chinese organizations such as (5) Baichuan-3~\footnote{\url{https://www.baichuan-ai.com}}, and (6) the newest GLM-4~\footnote{\url{https://open.bigmodel.cn/dev/api\#glm-4}} from the GLM series \citep{glm}.
\end{itemize}

The models listed serve general purposes and, for the most part, have been subjected to multiple stages of alignment refinement.

\item \textbf{Open-source models} encompass: (1) general-purpose architectures such as DeepSeek-LLM-Chat 67B \citep{deepseek-llm}, Qwen 72B \citep{qwen}, SeaLLM-v2 7B \citep{seallm}, and ChatGLM3 6B \citep{chatglm3}; (2) as well as models optimized for mathematics, including InternLM2-Math 20B~\footnote{\url{https://github.com/InternLM/InternLM-Math}}, which extends InternLM2 by further training on mathematics before undergoing instruction tuning; (3) Math-Shepherd-Mistral 7B, which leverages PPO training \citep{schulman2017proximal} atop Mistral 7B \citep{mistral} using a process-supervised reward mechanism; (4) the WizardMath family \citep{wizardmath}, which enhances mathematical capabilities in Mistral 7B and Llama-2 70B \citep{llama2} via evolve-instruct—a variant of instruction tuning reliant on AI-generated prompts—alongside PPO training, with most training data sourced from GSM8K and MATH; (5) MetaMath 70B \citep{metamath}, realized by fine-tuning Llama-2 70B on an enriched version of GSM8K and MATH; (6) ToRA 34B \cite{tora}, which adapts CodeLlama 34B for tool-supported mathematical reasoning through fine-tuning; and (7) MAmmoTH 70B \citep{MathInstruct}, a Llama-2 70B variant instruction-tuned with MathInstruct.

Table \ref{tab:sft_rl_math} highlights that, in evaluation scenarios where external tool utilization is restricted, DeepSeekMath-Instruct 7B attains notable stepwise reasoning proficiency. In particular, on the competition-level MATH benchmark, the model consistently outperforms all available open-source systems and outstrips most proprietary solutions—such as Inflection-2 and Gemini Pro—by at least 9\% in absolute terms. This performance edge holds even compared to significantly larger models (e.g., Qwen 72B) and those reinforced using math-specific RL (e.g., WizardMath-v1.1 7B). Although DeepSeekMath-Instruct is comparable to Chinese proprietary solutions like GLM-4 and Baichuan-3 on the MATH dataset, it remains behind GPT-4 and Gemini Ultra.

When the evaluation permits both natural language reasoning and code/tool-assisted strategies for problem solving, DeepSeekMath-Instruct 7B approaches 60\% accuracy on the MATH benchmark, overtaking all open-source competitors. On additional test suites, its performance aligns with that of DeepSeek-LLM-Chat 67B, the previous state-of-the-art open model despite having one-tenth the parameter count.

\section{Reinforcement Learning}

\subsection{Group Relative Policy Optimization} 
Reinforcement learning (RL) has shown considerable effectiveness in further refining LLMs' mathematical reasoning following the Supervised Fine-Tuning (SFT) phase \citep{wang2023math,wizardmath}. Here, we introduce a new RL method—Group Relative Policy Optimization (GRPO)—that is both resource-efficient and effective.

\subsubsection{From PPO to GRPO}
Proximal Policy Optimization (PPO) \citep{schulman2017proximal} serves as a widely adopted actor-critic RL framework for fine-tuning LLMs via reinforcement learning \citep{ouyang2022training}. PPO optimizes policy networks by maximizing a clipped surrogate objective, expressed as:

\begin{equation}
\footnotesize
    \mathcal{J}_{PPO}(\theta) = \mathbb{E}{[q \sim P(Q), o \sim \pi_{\theta_{old}}(O|q)]} \frac{1}{|o|} \sum_{t=1}^{|o|} \min \left[ \frac{\pi_\theta(o_{t} | q, o_{<t})}{\pi_{\theta_{old}}(o_{t} | q, o_{<t})} A_{t}, \text{clip} \left( \frac{\pi_\theta(o_{t} | q, o_{<t})}{\pi_{\theta_{old}}(o_{t} | q, o_{<t})}, 1 - \varepsilon, 1 + \varepsilon \

Here, $\pi_{\theta}$ denotes the current policy, while $\pi_{\theta_{old}}$ stands for the policy prior to the most recent update. The variables $q$ and $o$ represent questions drawn from a dataset and the corresponding responses sampled using the old policy $\pi_{\theta_{old}}$. The hyper-parameter $\varepsilon$ is related to the clipping mechanism in PPO, serving to enhance the robustness of the learning process. The term $A_t$ signifies the advantage, derived through the application of Generalized Advantage Estimation (GAE) \citep{gae}, utilizing future rewards $\{r_{\ge t}\}$ and a learned value function $V_{\psi}$. Consequently, PPO necessitates the concurrent training of both a value function and a policy. To limit excessive optimization of the reward model, the common technique is to append a token-level KL divergence penalty—calculated with respect to a reference policy—at every time step within the reward, as proposed in \citep{ouyang2022training}:

\begin{equation}
    r_{t} = r_\varphi(q, o_{\le t}) - \beta \log\frac{\pi_{\theta}(o_{t}|q, o_{<t})}{\pi_{ref}(o_{t}|q, o_{<t})},
\label{eq:PPO-reward}
\end{equation}

In this expression, $r_\varphi$ denotes the output of the reward model, $\pi_{ref}$ refers to the reference policy—often initialized from the SFT model—and $\beta$ controls the influence of the KL term.

Since the value function in PPO is typically realized by a neural model of a similar scale to the policy itself, this dual-model approach substantially increases memory usage and computational overhead. Furthermore, within RL-based fine-tuning, the value function is employed as a variance-reducing baseline when estimating the advantage. Notably, for LLMs, reward signals are commonly allocated only to the final token by the reward model, potentially complicating the accurate token-wise value prediction. To overcome these limitations, as depicted in Figure \ref{fig:grpo}, we introduce Group Relative Policy Optimization (GRPO), which eliminates the need for a separate value function as in PPO. Instead, GRPO utilizes the average reward of a collection of sampled outputs corresponding to the same prompt as a baseline. Concretely, given a question $q$, a group of outputs $\{o_1, o_2, \cdots, o_G\}$ is drawn from $\pi_{\theta_{old}}$, and the optimization target for the updated policy is given by:

\begin{equation}
\footnotesize
\begin{split}
    \mathcal{J}_{GRPO}(\theta) &= \mathbb{E}{[q \sim P(Q), \{o_i\}_{i=1}^G \sim \pi_{\theta_{old}}(O|q)]}  \\
    & \frac{1}{G}\sum_{i=1}^G\frac{1}{|o_i|} \sum_{t=1}^{|o_i|} \left\{ \min \left[ \frac{\pi_\theta(o_{i,t} | q, o_{i,<t})}{\pi_{\theta_{old}}(o_{i,t} | q, o_{i,<t})} \hat{A}_{i,t}, \text{clip} \left( \frac{\pi_\theta(o_{i,t} | q, o_{i,<t})}{\pi_{\theta_{old}}(o_{i,t} | q, o_{i,<t})}, 1 - \varepsilon, 1 + \varepsilon \right)  \hat{A}_{i,t} \right] - \beta \mathbb{D}_{KL}\left[\pi_{\theta} || \pi_{ref}\right]\right\} ,
\end{split}
\label{eq:GRPO-obj}
\end{equation}

The hyper-parameters $\varepsilon$ and $\beta$ serve as in the original formulation, and $\hat{A}_{i,t}$ denotes the advantage, computed with respect to the relative performance of outputs within the same group; details on its calculation will be elaborated in subsequent sections. GRPO's group-based approach to advantage estimation harmonizes with the inherently comparative design of reward models, which are often trained using pairwise output preferences conditioned on the same prompt. Furthermore, GRPO departs from PPO by adding the KL regularization term directly to the loss function instead of embedding it in the reward, thereby simplifying the advantage computation for $\hat{A}_{i,t}$. Unlike the reward-based KL penalty used in (\ref

\begin{algorithm}[t]
  \small
  \caption{Iterative Group Relative Policy Optimization}
  \textbf{Input} Initial policy $\pi_{\theta_{\text{init}}}$; reward model(s) $r_\varphi$; task prompt set $\mathcal{D}$;
  hyperparameters $\varepsilon$, $\beta$, $\mu$
  \begin{algorithmic}[1]
    \State Set policy $\pi_\theta \leftarrow \pi_{\theta_{\text{init}}}$
    \For{iteration = 1 to $I$}
      \State Set the reference policy $\pi_{ref} \leftarrow \pi_{\theta}$
      \For{step = 1 to $M$}
      \State Draw a batch $\mathcal{D}_b$ from $\mathcal{D}$
      \State Copy the current policy $\pi_{\theta_{old}} \leftarrow \pi_{\theta}$
      \State For each query $q$ in $\mathcal{D}_b$, generate $G$ outputs $\{o_i\}_{i=1}^G \sim \pi_{\theta_{old}} (\cdot \mid q)$
      \State For all sampled outputs $o_i$, obtain rewards $\{r_i\}_{i=1}^{G}$ via $r_\varphi$
      \State Estimate advantages $\hat{A}_{i,t}$ for the $t$-th token of $o_i$ using group-relative advantage
      \For{GRPO iteration = 1 to $\mu$}
        \State Update $\pi_{\theta}$ to maximize the GRPO objective in Equation \ref{eq:GC-GRPO}
      \EndFor
    \EndFor
    \State Periodically update $r_\varphi$ through replay-based continual learning.
    \EndFor
  \end{algorithmic}
  \textbf{Output} $\pi_\theta$
  \label{alg:iter-grpo}
\end{algorithm}

\subsubsection{Outcome Supervision RL with GRPO}
Let each prompt $q$ be associated with a set of outputs $\{o_1, o_2, \cdots, o_G\}$, generated by sampling from the old policy $\pi_{\theta_{old}}$. These outputs are evaluated by a reward model, resulting in reward values $\mathbf{r} = \{r_1, r_2, \cdots, r_G\}$. The rewards are standardized by first subtracting the mean of $\mathbf{r}$ and then dividing by its standard deviation. Within the outcome supervision framework, the normalized score at the end of output $o_i$ is assigned as the advantage $\hat{A}_{i, t}$ for every token $t$ in that output: specifically, $\hat{A}_{i, t} = \widetilde{r}_i = \frac{r_i- {\rm mean}(\mathbf{r})}{{\rm std}(\mathbf{r})}$. Policy improvement is performed by maximizing the objective in Equation (\ref{eq:GRPO-obj}).

\subsubsection{Process Supervision RL with GRPO}
Outcome-based reward signals, given only after a full output sequence, may lack sufficient granularity to train policies on complex math problems efficiently. To address this, we implement process supervision as described by \cite{wang2023math}, where feedback is given at each reasoning step rather than just the final answer. Given $q$ and $G$ sampled completions $\{o_1, o_2, \cdots, o_G\}$, a process-level reward model computes a reward at each step for every sampled output, yielding $\mathbf{R} = \{ \{r_1^{index(1)},\cdots,r_1^{index(K_1)}\}, \cdots,  \{r_G^{index(1)},\cdots,r_G^{index(K_G)}\} \}$, with $index(j)$ marking the token position at step $j$ and $K_i$ as the number of steps for output $i$. As with outcome rewards, stepwise rewards are standardized: $\widetilde{r}_i^{index(j)} = \frac{r_i^{index(j)} - {\rm mean}(\mathbf{R})}{{\rm std}(\mathbf{R})}$. The advantage $\hat{A}_{i, t}$ for token $t$ is computed by summing all future normalized stepwise rewards, that is, $\hat{A}_{i, t} = \sum_{index(j) \ge t} \widetilde{r}_i^{index(j)}$. The policy model is then trained to optimize the same objective

\subsubsection{Iterative RL with GRPO} As reinforcement learning advances, the existing reward model may no longer effectively supervise the evolving policy model. To address this, we investigate an iterative RL approach using GRPO. In Algorithm \ref{alg:iter-grpo}, this method involves generating updated reward model training datasets by sampling outputs from the current policy model. The prior reward model is updated continually through a replay strategy that integrates 10\% of previous data. The reference model is then synchronized with the policy model, after which the policy model undergoes continual training with the refreshed reward model.

\subsection{Training and Evaluating DeepSeekMath-RL}

Our RL experiments utilize DeepSeekMath-Instruct 7B as the base policy model. The RL training set is constructed from chain-of-thought questions pertaining to GSM8K and MATH within the SFT dataset, totaling approximately 144K items. We intentionally omit other SFT-derived questions to analyze RL's effectiveness on evaluation benchmarks where related data is withheld during RL. Following the methodology in \citep{wang2023math}, the reward model training set is compiled. The initial reward model is trained with DeepSeekMath-Base 7B at a learning rate of 2e-5. During GRPO-based training, the policy model adopts a learning rate of 1e-6, and the KL coefficient is fixed at 0.04. For each question, $64$ completions are sampled. Maximum sequence length is 1024 and training batch size is likewise set at 1024. The policy model is updated only once after each phase of exploration. We benchmark DeepSeekMath-RL 7B using the same protocol as DeepSeekMath-Instruct 7B. In this context, GSM8K and MATH tasks utilizing chain-of-thought reasoning are classified as in-domain, while all other evaluated benchmarks are considered out-of-domain.

Table \ref{tab:sft_rl_math} presents a comparative analysis of both open- and closed-source models, evaluating their capabilities in chain-of-thought and tool-augmented reasoning on English and Chinese tasks. The results indicate: 1) When employing chain-of-thought reasoning, DeepSeekMath-RL 7B achieves 88.2\% accuracy on GSM8K and 51.7\% on MATH. These results outperform all open-source competitors within the 7B to 70B parameter range, and even exceed most closed-source systems. 2) Notably, DeepSeekMath-RL 7B’s training is exclusively grounded in chain-of-thought-format instruction tuning data derived from GSM8K and MATH, starting from DeepSeekMath-Instruct 7B as the initialization. Despite this limited training data scope, DeepSeekMath-RL 7B consistently exceeds the performance of DeepSeekMath-Instruct 7B across all evaluation measures, highlighting the substantial advantages conferred by reinforcement learning.

\section{Discussion}
This section elaborates on insights gained from our investigations into pre-training and reinforcement learning.

\subsection{Lessons Learnt in Pre-Training}

We begin by sharing observations from the pre-training stage. Unless stated otherwise, the experiments follow the training configurations described in Section \ref{sec:quality-policy}. In this context, the DeepSeekMath Corpus refers to the dataset with 89B tokens accumulated from the second round of data collection.

\subsubsection{Code Training Benefits Mathematical Reasoning}

There is a prevailing but insufficiently tested belief that training models on code enhances their reasoning proficiency. Our findings offer empirical support for this claim within the mathematical domain: incorporating code in the training regimen elevates the model's mathematical reasoning abilities, regardless of whether external tools are used.

To explore the relationship between code training and mathematical reasoning, we devised experiments that incorporate both two-stage and one-stage training strategies:

\textbf{Two-Stage Training}

\begin{itemize}[topsep=0pt]
    \item \textbf{Code Training for 400B Tokens $\rightarrow$ Math Training for 150B Tokens}: DeepSeek-LLM 1.3B is initially exposed to 400B tokens derived from code sources, after which it receives a further 150B tokens focused on mathematical content;
    \item \textbf{General Training for 400B Tokens $\rightarrow$ Math Training for 150B Tokens}: For comparison, we also implement a regimen where, in the first stage, the model is trained on 400B tokens sampled from a broad, general-purpose corpus constructed by DeepSeek-AI, as opposed to code tokens. This approach enables examination of whether pretraining on code offers distinct benefits for mathematical reasoning relative to more generic textual data.
\end{itemize}

\textbf{One-Stage Training}

\begin{itemize}[topsep=0pt]
    \item \textbf{Math Training for 150B Tokens}: DeepSeek-LLM 1.3B undergoes training on 150B mathematical tokens exclusively;
    \item \textbf{Training on a mixture of 400B Code Tokens and 150B Math Tokens}: Since supplementing math training after code pretraining appears to reduce proficiency on code-related tasks, we also assess whether integrating both code and math tokens within a unified, single-phase training paradigm could simultaneously reinforce mathematical reasoning abilities and address the concern of catastrophic forgetting.
\end{itemize}

\paragraph{Results}
Tables \ref{tab:code-math} and \ref{tab:code-math-general-eval} present the comparative performance outcomes of the various training configurations.

Pretraining on code significantly enhances program-assisted mathematical reasoning, whether models are trained sequentially in two stages or within a combined single-stage setup. Table \ref{tab:code-math} illustrates that code pretraining in a two-stage schedule already markedly boosts model capabilities in solving GSM8K and MATH questions using Python. Adding subsequent math-focused training further elevates this performance. Notably, in the single-stage approach, jointly training with both code and math tokens helps prevent the decline in coding ability often observed in two-stage pipelines, and further advances both coding (Table \ref{tab:code-math-general-eval}) and program-aided mathematical reasoning (Table \ref{tab:code-math}) outcomes.

Moreover, pretraining with code enhances mathematical reasoning even in settings where external tools are not employed. In the two-stage sequence, beginning with code tokens leads to clear gains and increases the effectiveness of later math-centric training, culminating in the most robust results. Conversely, merging code and math tokens during single-stage training appears to hinder tool-free mathematical reasoning. We hypothesize this limitation arises from DeepSeek-LLM 1.3B’s restricted model capacity, which constrains its ability to simultaneously internalize both programming and mathematical information in a single training phase.

\subsubsection{ArXiv Papers Appear to Offer Limited Gains in Mathematical Reasoning}

It is standard practice to incorporate arXiv papers into the pre-training datasets for mathematical language models \citep{minerva,gpt-f,llemma,mathpile}. Nevertheless, systematic investigations into how arXiv content influences mathematical reasoning abilities are lacking. Somewhat unexpectedly, our findings suggest that exposure to arXiv papers does not significantly enhance mathematical reasoning performance. We evaluated models at varying parameter scales—specifically DeepSeek-LLM 1.3B and DeepSeek-Coder-Base-v1.5 7B \citep{deepseek-coder}—utilizing several processed arXiv corpora:

\begin{itemize}[topsep=0pt]
    \item \textbf{MathPile} \citep{mathpile}: A dataset spanning 8.9B tokens, curated through multiple filtering and cleaning procedures, with scientific arXiv submissions comprising over 85\% of its content;
    \item \textbf{ArXiv-RedPajama} \citep{redpajama}: The full arXiv repository’s LaTeX sources, stripped of preambles, annotations, macros, and bibliography sections, resulting in 28.0B tokens.
\end{itemize}

In our assessments, DeepSeek-LLM 1.3B was pre-trained on 150B tokens, while DeepSeek-Coder-Base-v1.5 7B was exposed to 40B tokens for each respective arXiv dataset. Empirical results show that exclusive pre-training on arXiv papers does not yield noticeable advancements—and may even impair performance—across the full spectrum of mathematical benchmarks considered. These include problem-solving tasks such as GSM8K and MATH (refer to Table \ref{tab:arxiv-cot}), multiple-choice evaluations like MMLU-STEM (also Table \ref{tab:arxiv-cot}), as well as theorem-proving tasks such as miniF2F (Table \ref{tab:arxiv_atp}).

Despite these observations, several caveats must be acknowledged. Specifically, our study does not address:

\begin{itemize}[topsep=0pt]
    \item The influence of arXiv data on narrowly defined or omitted mathematical tasks, such as the informal restatement of formal theorems or proofs;
    \item Potential interactions between arXiv material and diverse dataset types when used in combination;
    \item Whether arXiv papers might be beneficial when deployed in substantially larger-scale models.
\end{itemize}

Accordingly, more comprehensive inquiries are necessary to clarify these issues, which we defer to future work.

\subsection{Perspectives on Reinforcement Learning}
\subsubsection{Toward a Unifying Framework} In this section, we propose a general framework to analyze a range of training approaches, including SFT, RFT, DPO, PPO, GRPO, and support our investigation with targeted experiments exploring the critical components of this unified paradigm. Formally, the gradient with respect to parameters $\theta$ for a given training approach is represented as:

\begin{equation}
    \nabla_{\theta}\mathcal{J}_{\textcolor{red}{\mathcal{A}}}(\theta) = \mathbb{E}[\underbrace{(q,o) \sim \textcolor{red}{\mathcal{D}}}_{Data \ Source}]\left( \frac{1}{|o|} \sum_{t=1}^{|o|}  \underbrace{GC_{{\mathcal{A}}}(q, o, t, \textcolor{red}{\pi_{{rf}}})}_{Gradient \ Coefficient}  \nabla_{\theta}\log \pi_{\theta}(o_t | q, o_{<t})\right).
\label{eq:objective}
\end{equation}

This formulation is structured around three principal components: 1) the \textit{Data Source $\mathcal{D}$}, which supplies the examples for model optimization; 2) the \textit{Reward Function $\pi_{{rf}}$}, responsible for producing the reward signal used during training; and 3) the \textit{Algorithm $\mathcal{A}$}, which translates both the training data and reward into the gradient coefficient $GC$, thereby determining the scale of the applied penalty or reward for each data instance. Within this general framework, we examine various representative approaches:

\begin{itemize}
    \item  \textbf{Supervised Fine-tuning (SFT)}: This method adapts the pretrained model to new tasks by further training it with human-curated SFT datasets.
    \item \textbf{Rejection Sampling Fine-tuning (RFT)}: In RFT, the SFT model undergoes an additional fine-tuning phase utilizing outputs generated by the SFT model itself, selectively retaining outputs that meet correctness criteria for the associated questions.
    \item \textbf{Direct Preference Optimization (DPO)}: DPO incrementally updates the SFT model through fine-tuning on enriched outputs generated by the SFT model, leveraging a pairwise loss based on output preferences.
    \item \textbf{Online Rejection Sampling Fine-tuning (Online RFT)}: Diverging from conventional RFT, Online RFT initiates with the SFT model as policy and continually refines this policy using augmented samples generated in real time by the current policy.
    \item \textbf{PPO/GRPO}: These techniques start by initializing the policy model with SFT weights, then further strengthen the policy via reinforcement learning on outputs produced by the active policy model.
\end{itemize}

The key elements of these methods are outlined in Table \ref{tab:data_policy}. For a comprehensive exposition of the derivation steps, consult Appendix \ref{app:analysis-rl}.

\paragraph{Insight on Data Source} We categorize data sourcing into two main types: online sampling and offline sampling. In online sampling, training data is generated through exploration by the policy model as it is being trained, whereas offline sampling utilizes data collected from the outputs of the initial SFT model. RFT and DPO are representative of offline methods, while Online RFT and GRPO embody the online approach.

As illustrated in Figure \ref{fig:rl-analysis}, our results reveal that Online RFT surpasses RFT by a large margin across two benchmark tasks. During the initial training phase, Online RFT performs similarly to RFT; however, as training progresses, Online RFT secures a marked lead, emphasizing the efficacy of online training paradigms. This observation aligns with intuition: early in training, the actor model closely mirrors the SFT model, resulting in training samples with only minor distinctions. As training advances, actor-generated samples diverge more significantly, and the real-time nature of online data acquisition proves increasingly beneficial.

\paragraph{Insight on Gradient Coefficient} The algorithms map the reward signal into a gradient coefficient to facilitate model parameter updates. We delineate two reward functions: `Rule' and `Model.' The Rule function determines response quality by checking answer correctness, while the Model function employs a reward model trained to assign scores, using data labeled by rule-based judgments. As detailed in Equations \ref{eq:GC-RFT} and \ref{eq:GC-GRPO}, GRPO and Online RFT diverge in how they process the gradient coefficient: GRPO dynamically adjusts the coefficient in accordance with reward values produced by the reward model, enabling nuanced reinforcement or penalty based on response merit. Conversely, Online RFT applies uniform positive reinforcement to all correctly answered responses and offers no penalization for incorrect ones, lacking this capacity for fine-grained feedback.

As illustrated in Figure \ref{fig:rl-analysis}, GRPO consistently outperforms online RFT, underscoring the effectiveness of adjusting the positive and negative gradient coefficients. Additionally, the results reveal that GRPO+PS achieves better outcomes than GRPO+OS, underscoring the advantages of adopting more granular, step-sensitive gradient adjustments. Our investigation also extends to iterative RL: specifically, we perform two iterations in our experimental setup. According to Figure \ref{fig:iter-rl}, iterative RL yields substantial performance improvements, particularly pronounced during the initial iteration.

\subsubsection{Why RL Works?} This study employs reinforcement learning using a subset of the instruction tuning dataset, which leads to marked performance gains over models relying solely on instruction tuning. To further elucidate the mechanisms by which reinforcement learning is effective, we analyze the Pass@K and Maj@K accuracies of both the Instruct and RL-enhanced models across two evaluation sets. Figure \ref{fig:maj-pass} illustrates that RL substantially raises Maj@K accuracy while having little impact on Pass@K. This pattern suggests that reinforcement learning fortifies the output distribution, thus enhancing the likelihood of producing a correct response among the TopK candidates; that is, \textbf{the primary benefit derives from elevating correct answers within the TopK selections rather than improving core model competency.} In a similar vein, \citep{wang2023making} report a \textbf{misalignment problem} encountered in reasoning tasks with SFT models, showing that reasoning outcomes in such models can be improved by leveraging a range of preference alignment methods \citep{yuan2023rrhf,song2023preference,wang2023making}.

\subsubsection{How to Achieve More Effective RL?}
\label{sec:effec_rl}
Our findings establish that RL is highly effective for mathematical reasoning challenges. We also put forward a generalized framework to reinterpret various prominent training approaches, characterizing all methods as either variants or abstractions of RL. According to Equation \ref{eq:objective}, this framework comprises three fundamental aspects: Data Source, Algorithm, and Reward Function. We outline prospective research avenues for each of these aspects.

\paragraph{Data Source} The choice of data source is foundational across training paradigms. In RL scenarios, the data source pertains to unlabeled question prompts alongside responses generated from the policy model. Our present experiments utilize only the prompts from the instruction tuning dataset and employ basic nucleus sampling for output generation. We hypothesize that this constraint might account for the RL pipeline’s observed improvements being mostly confined to Maj@K metrics. In subsequent work, we plan to test the RL pipeline on out-of-distribution questions and integrate \textbf{more sophisticated sampling (decoding) strategies}, such as those rooted in tree-search methods \citep{yao2023tree}. Furthermore, advancements in \textbf{efficient inference techniques} \citep{xia-etal-2023-speculative,leviathan2023fast,kwon2023efficient,xia2024unlocking}, which critically determine the exploration capacity of policy models, are expected to have a substantial impact.

\paragraph{Algorithms} Algorithms utilize both data and the reward signal to calculate gradient coefficients for model parameter updates. According to Equation \ref{eq:objective}, prevailing approaches fundamentally \textbf{TRUST} the reward function’s signal to adjust the conditional probability—either boosting or diminishing it—for a particular token. Nonetheless, the reliability of the reward signal cannot always be guaranteed, particularly when confronted with highly complex tasks. For instance, even within meticulously labeled resources such as the PRM800K dataset \citep{lightman2023let}, approximately 20\% of the annotations are erroneous, despite extensive training of the annotators\footnote{\url{https://github.com/openai/prm800k/issues/12\#issuecomment-1728491852}}. This motivates an investigation into reinforcement learning algorithms that can maintain robustness in the face of noisy reward signals. We anticipate that \textbf{WEAK-TO-STRONG} \citep{burns2023weak} alignment approaches could significantly reshape the design and functioning of learning algorithms.

\paragraph{Reward Function} The reward function acts as the fundamental provider of training feedback. In reinforcement learning, it is typically operationalized as a neural reward model. We identify three primary research avenues for advancing reward models: 1) \textbf{Strategies to boost the generalization capability of the reward model.} It is crucial that the reward model can generalize effectively to out-of-distribution queries and sophisticated decoding scenarios; if it cannot, reinforcement learning may simply anchor the LLM’s output distribution rather than driving real capability enhancements; 2) \textbf{Mechanisms to represent and utilize reward model uncertainty.} Capturing uncertainty has the potential to serve as an essential connector between weak reward signals and weak-to-strong algorithmic strategies; 3) \textbf{Techniques for constructing high-fidelity process reward models efficiently,} thereby delivering granular guidance throughout reasoning tasks \citep{lightman2023let,wang2023math}.

\section{Conclusion, Limitation, and Future Work} 

In this work, we introduce DeepSeekMath, a model that surpasses all currently available open-source counterparts on the challenging MATH benchmark, narrowing the gap to the best closed-source systems. DeepSeekMath~was initialized with DeepSeek-Coder-v1.5 7B and further trained continuously on 500B tokens, including a substantial subset of 120B mathematical tokens harvested from Common Crawl. Through comprehensive ablation experiments, we demonstrate that web-scraped pages are a valuable resource for high-quality mathematical content, whereas the contributions from arXiv data were less significant than anticipated. We propose a new method, Group Relative Policy Optimization (GRPO), which builds on Proximal Policy Optimization (PPO) and enhances mathematical reasoning abilities while lowering memory requirements. Empirical evidence indicates that GRPO maintains its effectiveness even when DeepSeekMath-Instruct 7B already attains high performance on benchmarks. Furthermore, we introduce a unified framework for interpreting a range of techniques and outline several promising avenues for developing more robust reinforcement learning methodologies.

Despite the high quantitative reasoning scores achieved by DeepSeekMath~, its proficiency with geometry and theorem-proving tasks lags behind that of proprietary models. In preliminary evaluations, the model struggled with questions concerning triangles and ellipses, possibly pointing to biases in pre-training and fine-tuning data selection. Additionally, DeepSeekMath~is limited by its parameter scale, resulting in lower few-shot learning capabilities compared to GPT-4. While GPT-4 benefits noticeably from few-shot prompts, DeepSeekMath~exhibits similar results under both zero-shot and few-shot settings. Looking forward, we aim to refine our data selection and engineering pipeline to assemble a more robust and diverse pre-training corpus. Moreover, we plan to investigate the potential directions discussed in Section \ref{sec:effec_rl} for advancing reinforcement learning approaches tailored to large language models.

\end{document}